\section{Roadmap Update after Experiment 14C}
\label{sec:roadmap-post-14c}

Experiment~14C resolves the main ambiguity left by the multiclass Experiment~14B. The weak-class transient is real, but its simplest event-specific explanation is false. Underlearned classes are not starved of output events. In fact, they typically generate more class-associated event traffic. Heterogeneous compute periods materially delay finite-horizon class-wise convergence when total local work is matched, but dense SGD and EF-TopK exhibit the same qualitative sensitivity. The event method adds a distinct systems effect: the nominal $p_iT_i$ compensation does not equalize the nonlinear first-passage traffic after thresholding and reset.

This distinction matters for the roadmap. It would be premature to add client-specific thresholds, class-dependent jump amplitudes, or an explicit fairness controller merely to equalize event counts. Such a method would optimize a quantity that Experiment~14C does not identify as the cause of weak-class performance. It could also destroy the spontaneous communication adaptation that has repeatedly been useful in the project.

\subsection{What remains open from the fairness audit}

Two questions should be retained as secondary research directions.

First, event-rate normalization remains a systems question even if it is not the immediate class-convergence solution. Slow clients can transmit thousands of coordinates per completion because the $T_i$ compensation enters before the nonlinear threshold. If a deployment has per-client bandwidth, energy, or packet-size constraints, one may later seek a period-aware threshold or jump normalization that controls communication fairness while preserving the desired expected optimization contribution. That future experiment should optimize an explicit communication objective rather than being presented as a cure for class starvation.

Second, the hybrid trajectory is unusually sensitive to very small floating-point perturbations near threshold surfaces. A paper-quality benchmark should therefore use a containerized numerical environment and report multiple seeds or hardware-stable aggregate statistics. This is a reproducibility requirement, not a reason to abandon the mechanism.

\subsection{Recommended next gate: compact CNN}

The main optimization mechanism has now survived
\begin{equation}
\text{quadratics}
\rightarrow
\text{convex logistic regression}
\rightarrow
\text{synthetic MLPs}
\rightarrow
\text{deep MLPs}
\rightarrow
\text{binary Fashion-MNIST}
\rightarrow
\text{ten-class Fashion-MNIST}.
\end{equation}
Experiment~14C does not expose an event-specific failure severe enough to justify another compensating mechanism before changing architecture. The natural next gate is therefore a compact convolutional network on Fashion-MNIST.

A suitable Experiment~15A should compare the existing event optimizer, EF-TopK, and dense asynchronous SGD on a small architecture such as
\begin{equation}
\mathrm{Conv}(1,16,3)
\rightarrow
\mathrm{ReLU/tanh}
\rightarrow
\mathrm{pool}
\rightarrow
\mathrm{Conv}(16,32,3)
\rightarrow
\mathrm{pool}
\rightarrow
\mathrm{FC}
\rightarrow10.
\end{equation}
The exact activation choice should be audited rather than inherited blindly from the MLP experiments. The first goal is not state-of-the-art Fashion-MNIST accuracy. It is to test whether temporal coordinate events remain effective when parameters are embedded in convolutional filters with strong weight sharing and spatially structured gradients.

Mandatory diagnostics should include event rate and parameter coverage separately for convolution kernels, convolution biases, and dense output tensors. If convolutional kernels show severe imbalance, a filter- or block-event representation becomes scientifically motivated. If the raw coordinate mechanism remains competitive, additional block machinery should again be deferred.

\subsection{Why clustering remains deferred}

Experiments~14B--14C still do not establish incompatible client optima or persistent latent client groups. A single global model can learn the full ten-class task even under severe label skew. The finite-horizon class problem is strongly influenced by asynchronous compute and class difficulty and is also visible in conventional optimizers. Clustering would therefore solve a stronger problem than the evidence currently requires.

A later clustered/personalized experiment should deliberately introduce multimodal tasks or domains, for example incompatible class subsets, client-specific image transformations, device domains, or different label semantics. Only then is it meaningful to ask whether event-sign agreement, firing rates, co-firing statistics, or event timing can serve as a low-communication similarity signal.

The updated main progression is
\begin{equation}
\boxed{
14\mathrm{A}\;\text{binary real-data PASS}
\rightarrow
14\mathrm{B}\;\text{multiclass PASS with class-transient limitation}
\rightarrow
14\mathrm{C}\;\text{compute/event-fairness mechanism audit}
\rightarrow
15\mathrm{A}\;\text{compact CNN architecture gate}.
}
\end{equation}
